%% ---------------------------------------------------------------------
\section{Para llevar}
\label{sec:ps-cierre}
%% ---------------------------------------------------------------------

\begin{parallevar}
\begin{enumerate}
  \item El controlador ve todo el lazo: válvula, proceso y transmisor. Se identifica moviendo la CO y midiendo la PV, con ganancias normalizadas a los rangos.
  \item Modelo FOPDT $\Kproc e^{-Ls}/(\tau s+1)$ con la tangente de máxima pendiente: $L$ en el corte con el valor inicial, $\tau$ hasta el $63$\,\% (o hasta el $100$\,\% en la variante de Ziegler y Nichols, que sobrestima $\tau$). En un primer orden la pendiente máxima está en $t=L$.
  \item Malla abierta: Ziegler-Nichols, Cohen-Coon (agresivas, buscan el decaimiento de un cuarto) y la regla empírica del coeficiente de ajustabilidad (conservadora). Todas dependen de $L/\tau$.
  \item Malla cerrada: $\Kcr$ y $\Pcr$ de la oscilación sostenida. Ziegler-Nichols (PI: $0.45\Kcr$, $\Pcr/1.2$) es oscilatoria; Tyreus-Luyben es más conservadora y \emph{no} es de un cuarto.
  \item Síntesis directa: $\Gc=\frac{1}{\Gp}\frac{T}{1-T}$. Con $T=1/(\lambda s+1)$, $\lambda$ es la constante de tiempo del lazo: menor $\lambda$, más rápido. Los ceros inestables y el retardo no se invierten: van en $T$.
  \item IMC: $Q=(\tilde G_p^-)^{-1}F$, estructura que retroalimenta solo el error del modelo, y $\Gc=Q/(1-\tilde G_pQ)$, que para modelos sencillos es un PID con un solo parámetro, $\lambda$. Revise siempre la estabilidad interna.
  \item Desempeño y robustez se oponen: las reglas agresivas dejan márgenes de ganancia cercanos a $1.3$ y se desestabilizan con cambios moderados del proceso.
\end{enumerate}
\end{parallevar}

\begin{profundiza}[Correcciones respecto a la versión anterior del capítulo]
(i) Se decía que en un primer orden la pendiente máxima coincide con el $63$\,\%; está al comienzo de la respuesta, en $t=L$. (ii) El PI de Ziegler-Nichols de malla cerrada es $0.45\Kcr$ y $\Pcr/1.2$, no $0.4\Kcr$ y $0.8\Pcr$. (iii) Tyreus-Luyben no es ``un método de caída de un cuarto'': es más conservador que Ziegler-Nichols a propósito. (iv) En síntesis directa se decía que con $\lambda$ pequeño ``el sistema no responde rápido''; es al revés: $\ts\approx4\lambda$. (v) En el ejemplo de Tyreus-Luyben el proceso aparecía como $\Gc$ y $\Pcr=6.2$ en lugar de $6.28$\,s; en el de fase no mínima, $\tau_2=2.5$ era $\Td=2.5$; en el IMC faltaba la $s$ en $U(s)=\frac{1}{\Kproc}(\tau s+1)R(s)$. (vi) Proceso inestable: la condición es $F(1/\tau_u)=1$, no $F(\tau_u)=1$. (vii) En la tabla IMC, la fila del primer orden con retardo mezclaba la ganancia del PID con el $\Ti$ del PI; se separó en una fila PI y otra PID, y se retiró la fila del subamortiguado con cero y retardo, cuyas fórmulas no se pudieron verificar. (viii) La ganancia de velocidad $K_v$ se define con la pendiente de la rampa. (ix) La tabla del coeficiente de ajustabilidad se conserva, marcada como regla empírica. (x) Se completaron las simulaciones pendientes y se calcularon todos los números. (xi) Notación: $\tau_i,\tau_d\to\Ti,\Td$; $\theta\to L$; $T$ del FOPDT $\to\tau$, para no confundirla con $T(s)$.
\end{profundiza}

%% ---------------------------------------------------------------------
\section{Problemas}
\label{sec:ps-problemas}
%% ---------------------------------------------------------------------

Los problemas marcados con \etiquetaM{} requieren \matlab/\simulink; los marcados con \etiquetaEstrella{} son de profundización.

\begin{table}[ht]
\centering
\caption{Respuesta experimental del proceso de temperatura del problema~\ref{prob:ps-ident}.}
\label{tab:ps-datos}
\small
\begin{tabular}{@{}l*{7}{c}@{}}
\toprule
$t$ (s) & 0 & 2 & 4 & 6 & 8 & 10 & 15\\
PV ($\grados$C) & 80.0 & 80.0 & 80.1 & 81.7 & 83.3 & 84.6 & 87.7\\
\midrule
$t$ (s) & 20 & 24 & 30 & 40 & 50 & 60 & 100\\
PV ($\grados$C) & 89.9 & 91.4 & 93.1 & 95.0 & 96.2 & 96.9 & 97.9\\
\bottomrule
\end{tabular}
\end{table}

\begin{problema}[identificación]\label{prob:ps-ident}
En un proceso de temperatura, la CO ($0$--$100$\,\%) pasa de $30$ a $45$\,\% y el transmisor está calibrado entre $50$ y $150\,\grados$C; los datos están en la \cref{tab:ps-datos}. (a) Halle $\Kproc$ normalizada con \eqref{eq:ps-ganancia} y en unidades físicas; ¿por qué coinciden aquí? (b) Con la tangente y el método del $63$\,\%, halle $L$ y $\tau$. (c) Repita con la variante de Ziegler y Nichols y compare $\tau$. (d) Calcule $r_c=L/\tau$ y ubique el proceso en la \cref{tab:ps-ajust}. (e) \etiquetaM{} Superponga el modelo $1.2e^{-4s}/(20s+1)$ sobre los datos y calcule el error cuadrático medio. (f) Derive $y(t)=\Kproc(1-e^{-t/\tau})$ y diga dónde está la pendiente máxima.
\end{problema}

\begin{problema}[reglas de malla abierta]\label{prob:ps-abierta}
Con $\Kproc=1.2$, $L=4$\,s y $\tau=20$\,s: (a) halle las sintonías P, PI y PID de Ziegler-Nichols; (b) las P, PI, PD y PID de Cohen-Coon; (c) la del coeficiente de ajustabilidad. (d) \etiquetaM{} Simule los tres PID con el retardo exacto (estructura estándar, $N=10$) y tabule $\OS$, $\ts$, IAE y $|u_c|_{\max}$ ante un escalón de referencia. (e) Ordénelos de más a menos agresivo. (f) \etiquetaM{} Repita (d) con un escalón de perturbación a la entrada del proceso. ¿Se mantiene el orden? ¿Por qué una regla puede ser mejor para seguir la referencia y peor para rechazar perturbaciones?
\end{problema}

\begin{problema}[malla cerrada]\label{prob:ps-cerrada}
Con $\Gp=1/((s+1)(2s+1)(4s+1))$ y control proporcional: (a) halle $\Kcr$, $\wcr$ y $\Pcr$ por sustitución $s=\jw$; (b) verifique con Routh y \etiquetaM{} simulando con $\Kc$ un poco por debajo, igual y un poco por encima de $\Kcr$; (c) sintonice un PID ideal con Ziegler-Nichols y con Tyreus-Luyben; (d) \etiquetaM{} simule y mida la relación de decaimiento; ¿cuál da el cuarto? (e) \etiquetaM{} multiplique la ganancia del proceso por $1.5$ y repita; ¿cuál tolera mejor el cambio?
\end{problema}

\begin{problema}[síntesis directa]\label{prob:ps-sintesis}
(a) Para $\Gp=2/(5s+1)$ y $T=1/(\lambda s+1)$, muestre que resulta un PI y halle $\Kc$ y $\Ti$. (b) Para $\Gp=3/((2s+1)(10s+1))$, muestre que resulta un PID ideal. (c) \etiquetaM{} Simule (b) con $\lambda=1$, $3$ y $5$\,s; reporte $\ts$ y $|u_c|_{\max}$ y enuncie la relación entre $\lambda$ y la velocidad. (d) Con $\Gp=(1-9s)/((15s+1)(3s+1))$, muestre que $T=1/(\lambda s+1)$ da un controlador inestable y que $T=(1-9s)/(\lambda s+1)^2$ no. (e) Lleve el último a la forma PID en serie con filtro y verifique para $\lambda=5$. (f) Explique por qué el cero inestable debe ir en $T$ y qué pasa con el grado relativo si no.
\end{problema}

\begin{problema}[diseño IMC]\label{prob:ps-imc}
(a) Para $\tilde G_p=2/(5s+1)$, forme $Q$ con un filtro de primer orden y verifique que con modelo perfecto $Y=FR$. (b) Halle $\Gc$ con \eqref{eq:ps-imc-pid} y compare con el problema~\ref{prob:ps-sintesis}(a). (c) Repita para $\Kproc/((\tau_1s+1)(\tau_2s+1))$ y verifique la fila de la \cref{tab:ps-imc}. (d) Para $\tilde G_p=1.5(1-s)/((s+1)(4s+1))$, factorice de forma simple y pasa-todo, forme los dos $Q$ y compare sus órdenes. (e) \etiquetaM{} Simule los dos diseños con $\lambda=2$ ante referencia y ante una perturbación $1/(4s+1)$ a la salida; ¿cuál prefiere? (f) ¿Qué característica de la respuesta queda fijada por el cero inestable, haga lo que haga el controlador?
\end{problema}

\begin{problema}[robustez]\label{prob:ps-robustez}
Con las tres sintonías PID del problema~\ref{prob:ps-abierta}: (a) \etiquetaM{} calcule los márgenes de ganancia y de fase con el retardo exacto; (b) halle el retardo adicional $\Delta L$ que tolera cada una, en segundos y en porcentaje de $L$; (c) \etiquetaM{} varíe $\Kproc$ en $\pm30$\,\%, $\tau$ en $\pm30$\,\% y $L$ en $\pm50$\,\% (cinco valores de cada uno) y cuente los lazos estables; (d) identifique el peor caso y confírmelo por simulación; (e) ordene las sintonías por robustez y compare con el orden por agresividad; (f) \etiquetaEstrella{} halle el factor por el que hay que multiplicar el $\Kc$ de la menos robusta para que tolere toda la malla, y diga cuánto desempeño nominal cuesta.
\end{problema}

\begin{problema}[IMC-PID con tiempo muerto]\label{prob:ps-imc-fopdt}
Con $\Kproc=1.2$, $\tau=20$\,s y $L=4$\,s: (a) aproxime el retardo por Padé de primer orden, factorice y forme $Q$; (b) obtenga $\Gc$ y deduzca \eqref{eq:ps-imc-fopdt}; (c) compare con la fila FOPDT-PI de la \cref{tab:ps-imc}: ¿de qué aproximación sale cada una? (d) \etiquetaM{} Evalúe $\lambda=3.2$, $6.7$ y $10$\,s con el retardo exacto y tabule $\OS$, $\ts$, IAE y $|u_c|_{\max}$. (e) \etiquetaM{} Halle el menor $\lambda$ estable y compárelo con $0.8L$. (f) Compare la mejor sintonía IMC con las del problema~\ref{prob:ps-abierta} en IAE y margen de fase.
\end{problema}

\begin{problema}[estabilidad interna y errores de modelado]\label{prob:ps-interna}
(a) Para el ejemplo~\ref{ej:ps-bequette}, deduzca la ecuación característica y la cota de $\lambda$ por Routh. (b) \etiquetaM{} Simule por debajo, en y por encima de la cota. (c) \etiquetaM{} Con $\Gp=(1-9s)/((15s+1)(3s+1))$ y $\lambda=5$, simule $Q_1$ y $Q_2$ y grafique $y$ y $u$ durante $200$\,s. (d) Halle el polo responsable de la divergencia, su constante de tiempo y el instante en que $|u|$ supera $50$; explique por qué $y$ no lo revela. (e) Defina estabilidad interna. (f) Si el proceso es $K/(5s+1)$ y el modelo $2/(5s+1)$, halle la ecuación característica del lazo IMC y el rango de $K$ estable. ¿Cuánto error de ganancia tolera?
\end{problema}

%% ---------------------------------------------------------------------
\section{Soluciones seleccionadas}
%% ---------------------------------------------------------------------

\paragraph{Problema~\ref{prob:ps-ident}}
(a) $\Delta y\approx17.9\,\grados$C: $\Kproc=(17.9/100)/(15/100)=1.19\approx1.2$, y en unidades físicas $1.19\,\grados$C/\%; coinciden porque el rango de la PV es de $100\,\grados$C. (b) La pendiente máxima de los datos, $0.8\,\grados$C/s entre $4$ y $6$\,s, da $L=3.9$\,s; el $63.2$\,\% ($91.3\,\grados$C) se alcanza en $t\approx23.8$\,s, así que $\tau_{63}\approx19.9$\,s. (c) $\tau_{ZN}=17.9/0.8=22.4$\,s, mayor. (d) $r_c\approx0.2$: justo en el límite entre la fila PI y la PID; se usa la PID. (e) ECM $=0.0019\,(\grados\text{C})^2$. (f) La pendiente $\frac{\Kproc}{\tau}e^{-t/\tau}$ es máxima en $t=0$; en $t=\tau$ vale el $37$\,\% de ese máximo.

\paragraph{Problema~\ref{prob:ps-abierta}}
(a) P: $\Kc=4.17$; PI: $3.75$, $13.3$\,s; PID: $5$, $8$\,s, $2$\,s. (b) P: $4.44$; PI: $3.82$, $9.42$\,s; PD: $5.35$, $\Td=0.99$\,s; PID: $5.76$, $9.10$\,s, $1.40$\,s. (c) $2.29$, $22$\,s, $1.82$\,s. (d) ZN: $\OS=87$\,\%, $\ts=42$\,s, IAE $9.7$, $|u_c|_{\max}=54$; CC: $95$\,\%, $48$\,s, $11.5$, $61$; ajustabilidad: $0$\,\%, $26$\,s, $8.0$, $25$. (e) CC, ZN, ajustabilidad. (f) Picos $0.229$, $0.231$ y $0.288$; IAE $1.90$, $1.89$ y $9.6$: el orden se invierte.

\paragraph{Problema~\ref{prob:ps-cerrada}}
(a) $8s^3+14s^2+7s+1+\Kc=0$: $\wcr=\sqrt{7/8}=0.935$\,rad/s, $\Kcr=11.25$, $\Pcr=6.72$\,s. (c) ZN: $6.75$, $3.36$\,s, $0.84$\,s; TL: $5.11$, $14.8$\,s, $1.07$\,s. (d) ZN: $\OS=48$\,\%, decaimiento $0.16$; TL: $\OS=11$\,\%, sin segundo pico. (e) Con la ganancia $\times1.5$: $56$\,\% y $23$\,\%.

\paragraph{Problema~\ref{prob:ps-sintesis}}
(a) $\Kc=5/(2\lambda)$, $\Ti=5$\,s. (b) $\Kc=4/\lambda$, $\Ti=12$\,s, $\Td=5/3$\,s. (c) $\ts=3.9$, $12.1$ y $19.8$\,s; $|u_c|_{\max}=43.8$, $14.6$ y $8.8$: menor $\lambda$, más rápido. (d) Polo en $+1/9$. (e) $\Kc=18/(2\lambda+9)$, $\Ti=18$\,s, $\Td=2.5$\,s, $\tau_f=\lambda^2/(2\lambda+9)$; con $\lambda=5$, $0.947$ y $1.32$\,s.

\paragraph{Problema~\ref{prob:ps-imc}}
(d) Simple: $\tilde G_p^+=1-s$, $Q=(s+1)(4s+1)/(1.5(\lambda s+1)^2)$, de orden $2$. Pasa-todo: $\tilde G_p^+=(1-s)/(1+s)$, $\tilde G_p^-=1.5/(4s+1)$, $Q=(4s+1)/(1.5(\lambda s+1))$, de orden $1$. (e) Con $\lambda=2$: $\ts=12.5$ y $10.0$\,s, IAE $5$ y $4$ ($2\lambda+1$ y $\lambda+2$); ante la perturbación, picos $0.52$ y $0.47$. Se prefiere la pasa-todo. (f) La respuesta inversa.

\paragraph{Problema~\ref{prob:ps-robustez}}
(a) GM $1.33$, $1.31$, $2.91$; PM $39^\circ$, $27^\circ$, $75^\circ$. (b) $\Delta L=2.2$\,s ($54$\,\%), $1.3$\,s ($33$\,\%) y $10.1$\,s ($252$\,\%). (c) $98$, $90$ y $125$ de $125$. (d) $1.3\Kproc$, $0.7\tau$, $1.5L$. (f) CC con $\Kc\times0.48=2.77$: $\OS=24$\,\%, $\ts=31$\,s, IAE $9.4$, mejor que la original.

\paragraph{Problema~\ref{prob:ps-imc-fopdt}}
(c) La fila PI sale de $e^{-Ls}\approx1-Ls$; la PID, de Padé. La versión anterior de la tabla combinaba la ganancia de la segunda con el $\Ti$ de la primera. (d) $\Kc=3.53$, $2.11$, $1.53$; $\OS=16$\,\%, $0$, $0$; $\ts=16.5$, $28.3$, $41.1$\,s; IAE $5.8$, $8.7$, $12.0$; $|u_c|_{\max}=38$, $23$, $16$. (e) $\lambda_{\min}\approx0.75$\,s frente a $3.2$\,s. (f) IMC con $\lambda=6.7$: IAE $8.7$ y PM $76^\circ$, frente a $9.7$ y $39^\circ$ de Ziegler-Nichols.

\paragraph{Problema~\ref{prob:ps-interna}}
(a) $\lambda>0.887$. (d) Polo en $+1/9$, constante de tiempo $9$\,s; $|u|>50$ en $t=27.95$\,s. El cero del proceso cancela el polo en la función $Y/R$. (f) $1+Q(\Gp-\tilde G_p)=0$ se reduce a $2\lambda s+K=0$: estable para todo $K>0$. Con un error puro de ganancia, el IMC con modelo de primer orden tolera cualquier ganancia positiva.

\respuestaspausas
